\documentclass[
	12pt,
	openright,
	oneside,
	a4paper,
	english,
	french,
	spanish,
	brazil
]{abntex2}

\usepackage{lmodern}
\usepackage[T1]{fontenc}
\usepackage[utf8]{inputenc}
\usepackage{indentfirst}
\usepackage{color}
\usepackage{graphicx}
\usepackage{microtype}
\usepackage{lipsum}
\usepackage{hyperref}

\definecolor{blue}{RGB}{41,5,195}

\hypersetup{
    colorlinks=true,
    linkcolor=blue,
    citecolor=blue,
    urlcolor=blue
}

\titulo{QUANTIFLUX: METODOLOGIA DE ORQUESTRAÇÃO DE ENXAMES DE AGENTES DE IA POR ONDAS E PLANO DIREITOR DE ENGENHARIA}
\autor{Bruno M. Banho}
\local{Brasil}
\data{2026}
\orientador{Hermes Agent Framework}
\instituicao{
  Axio Engineering \& Meta-Cognitive Systems
  \par
  Laboratório de Arquitetura de Software e Agentes Autônomos
}
\tipotrabalho{Relatório Técnico de Engenharia}

\makeindex

\begin{document}

\frenchspacing

% CAPA
\imprimirfolhaderosto

% RESUMO
\begin{resumo}
 Este documento apresenta o \textbf{Quantiflux}, um arcabouço metanalítico e formal para orquestração de múltiplos agentes autônomos de inteligência artificial em projetos de engenharia de software de alta complexidade. A metodologia fundamenta-se na separação estrita de quatro camadas de abstração (Conceitual, Estratégico, Tático e Operacional) e na execução orientada a Ondas (\emph{Wave-based Execution}) derivadas de um Plano Diretor de Engenharia. Demonstra-se como o isolamento de contexto reduz drasticamente o \emph{context bleed} e previne a degradação conceitual do sistema durante a codificação massiva por agentes de linguagem.
 
 \vspace{\onelineskip}
 \noindent
 \textbf{Palavras-chave}: Orquestração de Agentes. Engenharia de Software. Plano Diretor. Execução por Ondas. Isolamento de Contexto.
\end{resumo}

% SUMÁRIO
\pdfbookmark[0]{\contentsname}{toc}
\tableofcontents*
\cleardoublepage

% CAPÍTULOS
\chapter{Introdução e Fundamentação}

\section{O Problema da Degradação Conceitual em Agentes Autônomos}
A utilização de modelos de linguagem de grande porte (LLMs) em tarefas de desenvolvimento de software trouxe capacidade de escrita de código em alta velocidade. Contudo, em projetos extensos, a atuação não estruturada de múltiplos agentes acarreta o fenômeno da \emph{crise conceitual} ou \emph{context drift}. Quando agentes especialistas atuam sem um plano unificado, decisões locais violam invariantes de arquitetura globais.

\section{A Filosofia Quantiflux}
O \textbf{Quantiflux} foi concebido para impor disciplina estrutural estrita antes da geração de qualquer linha de código. A premissa fundamental do framework é que a janela de contexto de um agente de nível estratégico é um ativo de alto valor. Para preservá-la, a execução de código é delegada a subagentes táticos temporários e isolados, operando sob o controle determinístico de um Plano Diretor de Engenharia (\emph{Master Engineering Plan}).

\section{Objetivos da Metodologia}
\begin{itemize}
    \item Formalizar a separação de responsabilidades em quatro níveis hierárquicos.
    \item Estabelecer a execução sequencial por Ondas (\emph{Waves}) com portões de qualidade invioláveis.
    \item Garantir que mudanças estruturais drásticas sejam tratadas no nível de tese e plano diretor, e não como improvisos operacionais.
\end{itemize}
\chapter{Arquitetura em Quatro Camadas}

\section{Estrutura de Abstração}
A orquestração Quantiflux divide o fluxo de trabalho em quatro níveis distintos e complementares:

\subsection{Nível Conceitual}
Representa a tese e a modelagem formal do sistema. Define as entidades fundamentais, os contratos imutáveis e as invariantes de domínio. Não contém detalhes de implementação ou bibliotecas específicas.

\subsection{Nível Estratégico}
Converte a tese conceitual no Plano Diretor de Engenharia (\emph{Master Engineering Plan}). Responsável por decompor o projeto em fases lógicas e ordenar as dependências entre módulos.

\subsection{Nível Tático}
Atua como o gerenciador de contexto e despachante. Filtra as informações relevantes do Plano Diretor para a Onda corrente e dispara subagentes especialistas via primitivas como \texttt{delegate\_task}.

\subsection{Nível Operacional}
Constituído pelos operadores e executores de código. Trabalham em ambientes limpos, escrevendo testes unitários e implementações locais. Submetem seus artefatos para validação nos portões de qualidade.
\chapter{Orquestração por Ondas e Portões de Qualidade}

\section{Decomposição por Ondas (Wave-based Execution)}
Em vez de permitir que agentes modifiquem múltiplos componentes simultaneamente, o Quantiflux exige a ordenação em Ondas:

\begin{enumerate}
    \item \textbf{Onda 0 (Core Contracts):} Definição de interfaces, DTOs e tipos abstratos.
    \item \textbf{Onda 1 (Adapters \& Persistência):} Implementação de adaptadores de infraestrutura e banco de dados.
    \item \textbf{Onda 2 (Serviços \& Lógica de Negócio):} Construção das regras de negócio consuming os contratos da Onda 0.
    \item \textbf{Onda N (Interface \& Integração Externa):} Camadas de exposição e APIs.
\end{enumerate}

\section{Isolamento de Contexto e Prevenção de Context Bleed}
Cada subagente despachado na Onda $K$ recebe estritamente as especificações da Onda $K$ e os contratos consolidados das Ondas $0 \dots K-1$. Informações de Ondas futuras são omitidas, eliminando alucinações e suposições prematuras.

\section{Portões de Qualidade (Wave Boundary Gates)}
A transição da Onda $K$ para a Onda $K+1$ é bloqueada até que os seguintes critérios sejam atendidos de forma auditável:
\begin{itemize}
    \item Compilação sem erros e aprovação de 100\% da suíte de testes unitários.
    \item Estado do repositório Git verificado e limpo (\texttt{git status}).
    \item Registro de métricas e lições aprendidas consolidado no repositório.
\end{itemize}
\chapter{Conclusão e Trabalhos Futuros}

\section{Considerações Finais}
O Quantiflux estabelece um padrão rigoroso de engenharia para o desenvolvimento assistido por enxames de agentes de IA. Ao tratar a estratégia de software como um processo determinístico alimentado por planos diretores e validado por portões de borda de onda, mitiga-se o risco de colapso arquitetural em projetos complexos.

\section{Evolução do Framework}
Como trabalhos futuros, prevê-se a integração nativa dos portões de onda a pipelines de CI/CD automatizados via GitHub Actions e o aprimoramento dos algoritmos de compressão de contexto para subagentes táticos.

\end{document}